\chapter{Corpus comparable}

\section{Définition et problèmatiques}

Les corpus comparables sont, comme définis par \citet{sharoff_2013a} des «collections de 
documents comparables dans leur contenu et leur forme à divers degrés et dimensions. cette définition inclue 
plusieurs types de corpus parallèles et non-parallèles, mais aussi des ensembles de corpus monolingues 
utilisés à but comparatif.» Par le passé plusieurs personnes ont étudier les corpus comparable et leur 
utilité comparé aux corpus parallèles.Les corpus parallèles sont très important dans la recherche liée au 
traitement automatique des langues. Cependant ces corpus peuvent poser plusieurs problèmes que précisent 
\citet{sharoff_2016} :
\begin{enumerate}
    \item Le nombre de paires languagières possibles augmente suivant le carré du nombre de langues.
    En utilisant des corpus parallèles, un bitext est nécessaire pour chaque paires languagières.
    En utilisant des corpus comparable, un corpus monolingue par langues suffit.
    \item Pour une traduction de meilleure qualité, des systèmes de traduction spécialisés sur des 
    genres et domaines particulier sont préférable. Mais il est bien plus difficile d'acquérir un
    corpus parallèles adaptés qu'un corpus comparable.
    \item  La langue évoluant au fil du temps, les corpus de formation doivent être mis à jour 
    régulièrement. Là encore, cela est plus difficile dans le cas d'un corpus parallèle.
  \end{enumerate}

Il est aussi, parfois, complexe de trouver des corpus parallèles dans les langues que l'on veut étudier.
Dans notre cas trouver des corpus parallèles pour des langues peu dotées et pratiquement impossible. Il est
donc primordial que l'on utilise des corpus comparable, plus abondant pour des langues peu dotées.

Les corpus existe à différents niveaux. Allant des textes parallèles aux textes sans relations et passant
par les textes fortement comparable et faiblement comparable \citet{sharoff_2013b}. Dans notre
cas il est intéressant de se pencher sur Wikipedia. Wikipedia possède des pages en Corse, Occitan, Alsacien et
Français qui possède l'avantage d'avoir des liens interlangues créés entre eux. Ces liens, créer par les auteurs
des articles, nous permettront de trouver des articles disponible dans les 4 langues étudiées. En l'occurence il 
est préférable pour nous de trouver des articles avec une forme textuel similaire. Par exemple un article sur
Paris serait un bon document pour créer un corpus comparable. Cet article étant disponible dans les 4 langues et 
ayant des similarités qui rendent le texte comparable.

Déjà dans l'état de l'art en 2016 on peux relever que des travaux sur les corpus comparable dans la traduction automatique
avait était réalisés. En effet des projets tel que Accurat\footnote{http://www.accurat-project.eu} ou 
TTC\footnote{https://www.ttc-project.eu} était en plein essort. Ces projets s'intéresse partiellement
ou totalement à la traduction automatique en utilisant des corpus comparables  \citep{sharoff_2016}.
Ces travaux nous montrent l'utilité d'étudier des corpus comparable et l'intérêt porter à cette recherche.

\newpage
\section{Fragments parallèles}

Blablabla

% 2016, à mtn
% explications sur les sentence transformers, bitext, 
% diff entre corpus parallèles et comparables

% Données plus de détails sur les recherches passées
% sous formes de sous-sections.

